\food{Mum's Chili Con Carne}
\preptime{20 mins} \cooktime{1 hour 20 mins} \people{6}

\recipe{
    \unit[750]{g} & Minced beef (10-15\% fat)\\
    \unit[6]{tbsp} & Oil\\
    2 & Onions (medium, finely chopped)\\
    \unit[2]{cloves} & Garlic (peeled and crushed)\\
    \unit[1-2]{tsp} & Chilli powder\\
    \unit[0.5-1]{tsp} & Cumin\\
    \unit[0.5-1]{tsp} & Dried oregano\\
    \unit[3-4]{tbsp} & Tomato pur\'ee\\
    \unit[0.5]{pint} & Beef stock\\
     & Salt\\
     & Black pepper (freshly ground)\\
    1 & 425g can red kidney beans (drained, or not)\\
     & Fresh oregano or chives, chopped, to garnish
}{
    \item Break up the mince with a fork and fry in 3 tablespoons of the oil until it begins to brown. Transfer to a casserole dish with a draining spoon.
    \item Fry the onions and garlic in the remaining 3 tablespoons of oil until just golden. Add the chilli powder (to taste), cumin and oregano, and cook over a gentle heat for a few minutes to bring out the flavours. Stir in the tomato pur\'ee.
    \item Add the onion mixture to the meat in the casserole. Add the stock, salt and pepper to taste. Cover and cook for around 1 hour, until the meat is tender.
    \item Skim the fat from the surface if you like, then add the kidney beans 15 minutes before serving.
    \item Garnish with chopped fresh oregano or chives, if using.
}
%\photo{chiliconcarne.jpg}

\info{Originally based on a recipe from Sallie Morris's \textit{Cooking with Herbs and Spices}.
If using dried kidney beans instead of canned: cover 175g of beans with cold water and leave to soak overnight. Drain, rinse and drain again, then place in a large pan with plenty of cold water. Bring to the boil and ensure the beans boil properly for the first 10 minutes of cooking, then reduce the heat, cover and simmer for 1--1.5 hours until tender, adding a little salt near the end. Drain and use as directed in the recipe.}
